% =====================================================================
\chapter{Introduction and the Equivalence Principle}
% =====================================================================

\section{Why General Relativity?}

Physics, as we know it today, describes Nature through \emph{four fundamental interactions}:
\begin{enumerate}[label=(\roman*)]
  \item \textbf{Electromagnetism}, described classically by Maxwell's equations and, at the quantum level, by Quantum Electrodynamics (QED);
  \item the \textbf{weak interaction}, responsible e.g.\ for $\beta$ decay;
  \item the \textbf{strong interaction}, described by Quantum Chromodynamics (QCD);
  \item \textbf{gravity}, described by Newtonian mechanics in the regime of weak fields and small velocities, and in full generality by \emph{General Relativity} (GR).
\end{enumerate}
The first two are unified, within the Standard Model (SM) of particle physics, into the \emph{electroweak theory},
based on the gauge group $SU(2)_L\times U(1)_Y$. Together with QCD (gauge group $SU(3)_c$) they form the SM,
which is a \emph{quantum} field theory. Gravity stands apart: General Relativity is a \emph{classical} (i.e.\ non-quantum) theory,
and no fully satisfactory quantum theory of gravity is known.

General Relativity is at the same time
\begin{itemize}[nosep]
  \item a \textbf{theory of gravity}, and
  \item a \textbf{theory of space-time} itself.
\end{itemize}
In Newtonian physics space and time are a fixed, absolute stage on which physics happens. In GR the stage becomes an actor:
the geometry of space-time is dynamical and is determined by its matter and energy content. This double role is beautifully summarised by J.\,A.~Wheeler:
\begin{quote}
\emph{``Space-time tells matter how to move; matter tells space-time how to curve.''}
\end{quote}
The first half of the sentence will be made precise by the \emph{geodesic equation} (\cref{ch:connection}), the second half by the \emph{Einstein field equations} (\cref{ch:efe}).

\section{Useful notions}

Before entering the subject we fix a few basic notions that will be used throughout the course.

\begin{definition}{Event}{event}
An \textbf{event} is a physical occurrence at a specific point of space and at a specific instant of time, i.e.\ associated with a specific space-time point.
An event has no spatial extension and no duration: it is a zero-dimensional object.
\end{definition}
\textit{Examples:} the emission of a photon by an atom; the presence of a point-like particle at a given place at a given time; two particles colliding.

\begin{definition}{Observer}{observer}
An \textbf{observer} is, essentially, a \emph{reference frame}: a way of assigning coordinates (three spatial coordinates and a time) to events.
\end{definition}
We can switch from one observer to another by means of \emph{coordinate transformations}. A central theme of the course is that physical statements must not depend on which observer (i.e.\ which coordinates) we use.

\begin{definition}{Local time}{localtime}
Each observer is equipped with an ideal standard clock and has a notion of time: he/she can assign a time to events and, therefore, a temporal ordering (``this event happened before that one'').
All observers measure time in the same way (with identical ideal clocks), even if, as we will see, they do not necessarily agree on the result.
\end{definition}

\begin{definition}{World-line}{worldline}
The \textbf{world-line} of a point-like object is the path that it follows in space-time. It is a one-dimensional object (a curve).
For extended objects one speaks of a \emph{world-sheet} (a string, 2D) or a \emph{world-volume} (higher-dimensional objects).
\end{definition}

% ---------------------------------------------------------------------
\section{The Equivalence Principle}
% ---------------------------------------------------------------------
The Equivalence Principle (EP) is the physical cornerstone on which Einstein built General Relativity. We first recall its Newtonian (Galilean) version and then formulate Einstein's version.

\subsection{The Newtonian (Galilean) version}

Consider a body of mass $m$, located at a point $P$ near the surface of the Earth. In Newtonian physics two \emph{a priori} completely different notions of mass enter its description:
\begin{itemize}
  \item From \textbf{Newtonian gravity}, the gravitational force acting on the body is
  \begin{equation}
    \vec F = -m_g\,\nabla\phi ,
  \end{equation}
  where $\phi$ is the gravitational potential and $m_g$ is the \textbf{gravitational mass}.
  It plays for gravity the role that the electric charge plays for electromagnetism ($\vec F=q\vec E=-q\nabla V$): it is a kind of ``gravitational charge'', measuring how strongly the body couples to the gravitational field.
  \item From \textbf{Newtonian mechanics} (Newton's second law),
  \begin{equation}
    \vec F = m_i\,\vec a ,
  \end{equation}
  where $m_i$ is the \textbf{inertial mass}. It tells us how much the body \emph{resists} being accelerated by \emph{any} force.
\end{itemize}

\begin{figure}[ht]
\centering
\begin{tikzpicture}[scale=1]
  \fill[pattern=north east lines] (-2,-0.25) rectangle (2,0);
  \draw[thick] (-2,0) -- (2,0) node[right] {Earth};
  \fill (0,1.6) circle (2.5pt) node[left=3pt] {$P$} node[right=3pt] {$m$};
  \draw[-{Stealth[length=3mm]},thick,red!70!black] (0,1.5) -- (0,0.55) node[midway,right] {$\vec F_g$};
\end{tikzpicture}
\caption{A body of mass $m$ at a point $P$ above the surface of the Earth, subject to the gravitational force $\vec F_g$.}
\label{fig:newton-ep}
\end{figure}

There is no reason, within Newtonian theory, for these two quantities to be equal: they are \emph{fundamentally different} concepts. Yet, experimentally, one observes (to an extremely high precision: $|m_i-m_g|/m_i\lesssim 10^{-15}$ in modern torsion-balance and space experiments) that
\begin{equation}
  m_i = m_g
\end{equation}
(more precisely, they are proportional to each other with a universal constant, which can be absorbed into the definition of Newton's constant $\GN$ and set to one).

\begin{keyresult}[Galilean Equivalence Principle]
The inertial and the gravitational masses of any body are equal, $m_i=m_g$. Consequently, equating the two expressions of the force,
\begin{equation}
  m_i\,\vec a = -m_g\,\nabla\phi \quad\Longrightarrow\quad \vec a = -\nabla\phi ,
\end{equation}
\emph{independently of the mass and of the composition of the body}. This is the \textbf{universality of free fall}: all bodies fall in the same way in a gravitational field (Galileo's experiments, the hammer and the feather on the Moon, \dots). In other words, \emph{the motion in a gravitational field does not depend on the mass}.
\end{keyresult}

This is a very peculiar property of gravity: for the electric force, $\vec a=(q/m)\vec E$ depends on the charge-to-mass ratio of the body, so different bodies fall differently in an electric field.

\subsection{Einstein's version of the Equivalence Principle}

Einstein's insight was to take the universality of free fall seriously and to turn it into a principle about \emph{reference frames}. Consider an elevator (a closed box) containing a particle of mass $m$ at a point $P$. We introduce two frames (see \cref{fig:elevator}):
\begin{itemize}[nosep]
  \item $R$, a ``static'' frame with coordinates $(x,z)$;
  \item $R'$, a frame \emph{comoving} with the elevator, with coordinates $(x',z')$.
\end{itemize}

\begin{figure}[ht]
\centering
\begin{tikzpicture}[scale=0.95,>=Stealth]
  % static frame R
  \draw[->,thick] (0,0) -- (7,0) node[right] {$x$};
  \draw[->,thick] (0,0) -- (0,4.4) node[above] {$z$};
  \node at (7.0,-0.45) {$(R)$ static};
  % comoving frame R'
  \draw[->] (1.2,1.0) -- (6.5,1.0) node[right] {$x'$};
  \draw[->] (1.2,1.0) -- (1.2,4.2) node[above] {$z'$};
  \node at (6.6,0.65) {$(R')$ comoving};
  % elevator
  \draw[very thick,fill=gray!10] (1.7,1.0) rectangle (3.4,3.4);
  \draw (2.55,3.4) -- (2.55,4.3); % rope
  \fill (2.55,2.5) circle (2.5pt) node[right=3pt] {$P$} node[above left=1pt] {$m$};
  \draw[->,thick,red!70!black] (2.55,2.4) -- (2.55,1.7) node[midway,left] {$\vec F_g$};
  \draw[->,very thick,blue!70!black] (4.0,2.2) -- (4.0,3.4) node[right] {$\vec g$};
\end{tikzpicture}
\caption{Einstein's elevator. $R$ is a static frame, $R'$ is the frame comoving with the elevator. Depending on the case, the elevator is at rest in empty space, uniformly accelerated (acceleration $\vec g$ upwards), or at rest in a uniform gravitational field (force $\vec F_g$ downwards on the particle).}
\label{fig:elevator}
\end{figure}

We consider the following situations.
\begin{description}
  \item[Case 0.] The elevator is in outer space, far from any mass, and no force acts on it. Then the particle simply \emph{floats} inside the elevator.
  \item[Case 1.] Still in outer space, but now the elevator is accelerated upwards with a constant acceleration $\vec g$ (e.g.\ by a rocket). Then the particle (which feels no force) is approached by the floor of the elevator: for an observer in the elevator, the particle ``falls'' towards the floor with acceleration $-\vec g$.
  \item[Case 2.] The elevator is at rest, but inside a constant gravitational field $\vec g$ pointing downwards (e.g.\ at rest on the surface of the Earth). Again, the particle approaches the floor with acceleration $-\vec g$.
\end{description}
Both $R$ and $R'$ agree on these statements. \textbf{But} the observer in $R'$ \emph{cannot distinguish case 1 from case 2} by any experiment performed inside the elevator. Hence:
\begin{center}
\fbox{\textsc{Gravity} $=$ \textsc{Acceleration}}
\end{center}
(locally, see below).

\paragraph{Quantitative version of case 1.} Let us neglect special-relativistic effects (so that $t'=t$). In $R$ the elevator floor moves as $z_{\rm floor}(t)=\tfrac12 g t^2$, so that $\ddot z_{\rm floor}=g$. The coordinates of the two frames are related by
\begin{equation}
  z' = z-\tfrac12 g t^2,\qquad x'=x,\qquad t'=t .
\end{equation}
Differentiating twice with respect to time, $\ddot z'=\ddot z-g$. In $R$, Newton's law reads $F=m\ddot z$. In $R'$,
\begin{equation}
  F' \equiv m\ddot z' = m\ddot z - mg \quad\Longrightarrow\quad F' = F - mg .
\end{equation}
The acceleration of $R'$ with respect to $R$ manifests itself in $R'$ as an additional force $-mg$, which is \emph{exactly} what a uniform gravitational field would produce. Note that this ``inertial force'' is proportional to the inertial mass, whereas gravity is proportional to the gravitational mass: their indistinguishability is precisely $m_i=m_g$.

\begin{description}
  \item[Case 3.] The elevator is in the gravitational field of the Earth and we cut the rope, so that it falls freely. Then the particle \emph{floats} with respect to $R'$, exactly as in case 0: the observer in $R'$ cannot tell whether he/she is at rest in outer space or freely falling in a gravitational field. Freely falling frames are locally equivalent to inertial frames: gravity has been ``switched off'' by going to a freely-falling frame.
\end{description}

\begin{keyresult}[Einstein's Equivalence Principle]
No physical experiment can distinguish \emph{locally} between a uniform gravitational field acting on a local inertial observer and a uniform acceleration with respect to that local inertial observer.

Here the \emph{local inertial observer} is the one for whom Special Relativity holds (a freely-falling observer).
\end{keyresult}

\subsection{Why ``locally''?}

Real gravitational fields are not uniform: the gravitational field of the Earth points towards its centre and its strength decreases with the distance. Consider a large freely-falling elevator (or two elevators at different positions, see \cref{fig:locally}). Two particles initially at rest at the two ends of a horizontal elevator both fall towards the centre of the Earth, along \emph{non-parallel} directions: seen from inside the elevator, they slowly approach each other. Similarly, two particles at different heights in a vertical elevator feel slightly different accelerations and drift apart. These \emph{tidal effects} cannot be removed by changing reference frame: they are the genuine signature of a gravitational field, and they will turn out to be the manifestation of space-time \emph{curvature} (\cref{sec:geodesic-deviation}).

\begin{figure}[ht]
\centering
\begin{tikzpicture}[scale=0.9,>=Stealth]
  \draw[thick,fill=brown!15] (-4.5,0) arc (160:20:4.79) ;
  \node at (0,0.4) {Earth};
  \coordinate (C) at (0,-1.6);
  % horizontal elevator (left)
  \begin{scope}[shift={(-2.6,3.5)},rotate=-15]
    \draw[thick] (-1.4,-0.35) rectangle (1.4,0.35);
    \fill (-1.1,0) circle (2pt); \fill (1.1,0) circle (2pt);
  \end{scope}
  \draw[->,red!70!black] (-3.65,3.75) -- ($(-3.65,3.75)!0.22!(C)$);
  \draw[->,red!70!black] (-1.55,3.2) -- ($(-1.55,3.2)!0.22!(C)$);
  % vertical elevator (right)
  \begin{scope}[shift={(2.5,4.0)},rotate=25]
    \draw[thick] (-0.35,-1.3) rectangle (0.35,1.3);
    \fill (0,1.0) circle (2pt); \fill (0,-1.0) circle (2pt);
  \end{scope}
  \draw[->,red!70!black] (2.08,4.9) -- ($(2.08,4.9)!0.12!(C)$);
  \draw[->,red!70!black] (2.92,3.1) -- ($(2.92,3.1)!0.2!(C)$);
  \draw[dashed,gray] (-3.65,3.75) -- (C);
  \draw[dashed,gray] (-1.55,3.2) -- (C);
\end{tikzpicture}
\caption{Why the Equivalence Principle only holds locally. In a freely-falling elevator, test particles fall towards the centre of the Earth along converging directions (left) and with height-dependent accelerations (right): these tidal effects cannot be removed by a change of frame.}
\label{fig:locally}
\end{figure}

\subsection{First consequence: light bends in a gravitational field}

Consider again the elevator accelerated upwards with acceleration $\vec g$ (case 1), and a laser fixed on the left wall that emits a light pulse horizontally (\cref{fig:lightbend}). In the inertial frame $R$ the light travels in a straight line. During the time the light needs to cross the elevator, however, the elevator has moved upwards: in $R'$ the light enters at the point $P_1$ and hits the opposite wall at a point $P_2$ which is \emph{lower} than $P_1$. Hence, in the accelerated frame, light follows a curved (parabolic) trajectory.

By the Equivalence Principle the same must happen in an elevator at rest in a gravitational field: \textbf{light feels gravity, and light bends}. This is a striking prediction, since photons are massless; it will be computed quantitatively in \cref{sec:light-bending}.

\begin{figure}[ht]
\centering
\begin{tikzpicture}[scale=0.95,>=Stealth]
  \draw[->,thick] (0,0) -- (6.5,0) node[right] {$x$};
  \draw[->,thick] (0,0) -- (0,4.2) node[above] {$z$};
  \node at (6.6,-0.4) {$(R)$};
  \draw[very thick,fill=gray!8] (1.5,0.8) rectangle (4.3,3.6);
  \draw[fill=black] (1.2,2.85) rectangle (1.5,3.05);
  \node[left] at (1.15,2.95) {laser};
  \fill (1.5,2.95) circle (1.8pt) node[above right] {$P_1$};
  \draw[thick,red!80!black,dashed] (1.5,2.95) .. controls (2.9,2.9) and (3.7,2.6) .. (4.3,2.25);
  \fill (4.3,2.25) circle (1.8pt) node[right] {$P_2$};
  \draw[gray,thin] (1.5,2.95) -- (4.3,2.95);
  \draw[->,very thick,blue!70!black] (5.2,2.0) -- (5.2,3.2) node[right] {$\vec g$};
  \node at (5.5,1.2) {$(R')$};
\end{tikzpicture}
\caption{A light ray emitted horizontally in an elevator accelerated upwards. In the comoving frame $R'$ the ray hits the opposite wall at $P_2$, lower than the emission height $P_1$: light bends. By the EP the same happens in a gravitational field.}
\label{fig:lightbend}
\end{figure}

\subsection{Second consequence: the gravitational redshift}\label{sec:redshift-ep}

Consider a tower of height $d$ in a uniform gravitational field $g$. A photon of frequency $f_P$ is emitted by an emitter $P$ at the bottom and detected by a detector $Q$ at the top. Which frequency $f_Q$ does $Q$ measure?

By the EP, we can replace the gravitational field by an elevator accelerated upwards with acceleration $g$, in which $P$ and $Q$ are at rest and separated by the distance $d$ (\cref{fig:redshift}). We work to first order in $gd/c^2\ll 1$.
\begin{itemize}
  \item The time of flight of the photon is $\Delta t= d/c$.
  \item During this time, the velocity of the elevator (hence of the detector $Q$ with respect to the inertial frame in which the emitter was at rest at the moment of emission) increases by
  \begin{equation}
    \Delta v = g\,\Delta t = \frac{g d}{c}.
  \end{equation}
  \item The detector is therefore moving \emph{away} from the point of emission with velocity $\Delta v$ when it receives the photon, and it observes a Doppler-shifted frequency (non-relativistic Doppler effect, since $\Delta v\ll c$):
  \begin{equation}
    \frac{f_Q}{f_P} = 1-\frac{\Delta v}{c} = 1-\frac{g d}{c^2}.
  \end{equation}
\end{itemize}

\begin{figure}[ht]
\centering
\begin{tikzpicture}[scale=0.9,>=Stealth]
  \draw[->,thick] (0,0) -- (7,0) node[right] {$x$};
  \draw[->,thick] (0,0) -- (0,4.6) node[above] {$z$};
  % at emission
  \draw[thick] (1.2,0.5) rectangle (2.0,1.3); \node[right] at (2.0,0.9) {$P$};
  \draw[thick] (1.2,3.0) rectangle (2.0,3.8); \node[right] at (2.0,3.4) {$Q$};
  \draw[<->] (0.8,0.9) -- (0.8,3.4) node[midway,left] {$d$};
  \draw[decorate,decoration={snake,amplitude=1.2pt,segment length=5pt},->,red!80!black] (1.6,1.35) -- (1.6,2.2) node[right] {$\gamma$};
  \node at (1.6,-0.4) {emission ($t$)};
  % at detection
  \draw[thick] (4.4,1.0) rectangle (5.2,1.8); \node[right] at (5.2,1.4) {$P$};
  \draw[thick] (4.4,3.5) rectangle (5.2,4.3); \node[right] at (5.2,3.9) {$Q$};
  \draw[decorate,decoration={snake,amplitude=1.2pt,segment length=5pt},->,red!80!black] (4.8,2.6) -- (4.8,3.45);
  \node at (4.8,-0.4) {detection ($t+d/c$)};
  \draw[->,very thick,blue!70!black] (6.4,1.6) -- (6.4,2.8) node[right] {$\vec g$};
\end{tikzpicture}
\caption{Gravitational redshift via the Equivalence Principle: emitter $P$ and detector $Q$, separated by $d$, are accelerated upwards with acceleration $g$. During the photon's time of flight $d/c$ the detector gains a velocity $\Delta v=gd/c$.}
\label{fig:redshift}
\end{figure}

\begin{keyresult}[Gravitational redshift]
By the Equivalence Principle (acceleration $\leftrightarrow$ gravity), in the presence of a uniform gravitational field and in a fixed reference frame, a photon climbing a height $d$ is redshifted:
\begin{equation}
  \frac{f_Q}{f_P} = 1-\frac{g d}{c^2} = 1-\frac{\Delta\phi}{c^2},\qquad \Delta\phi=\phi_Q-\phi_P = g d .
\end{equation}
\end{keyresult}

\paragraph{Interpretation in terms of clocks.} The number of wave crests is conserved between emission and reception (in a static situation nothing accumulates between $P$ and $Q$). The fact that $Q$ counts fewer crests per unit of its own time than $P$ emits per unit of its own time means that the two clocks do not tick at the same rate: equivalently, \textbf{clocks at lower gravitational potential tick slower}. A clock at the bottom ($P$), seen from the top ($Q$), runs slow by the factor $1-gd/c^2$ (and, conversely, the clock at $Q$ appears to run fast when seen from $P$). Time itself flows differently at different heights: gravity affects the geometry of \emph{space-time}, not just the motion of bodies. This effect has been measured in the Pound--Rebka experiment (1959) and is routinely taken into account by the GPS system.
